\chapter{Das Smart-Home im Unterricht: Konzeptionelle Lernszenarien}\label{chap:unterrichtsszenarien}
\todo[color=yellow]{Fünf konzeptionelle Unterrichtsmodule ausarbeiten.}

\section{Methodischer Zugang}\label{sec:methodischer-zugang}
\todo[color=yellow]{Forschendes und projektbasiertes Lernen beschreiben.}

\section{Modul 1: Strom und Elektronik im Smart-Home}\label{sec:modul-strom-elektronik}
\todo[color=yellow]{Stromkreise, Wasser-Modell und Relais vermitteln.}

\section{Modul 2: Smart-Home und Umwelt}\label{sec:modul-umwelt}
\todo[color=yellow]{Energieeffizienz, Sensorik und Vampirstrom behandeln.}

\section{Modul 3: Automatisierung und Logik}\label{sec:modul-automatisierung-logik}
\todo[color=yellow]{Bedingungen, Zustände und Requests erklären.}

\section{Modul 4: Netzwerke und IoT}\label{sec:modul-netzwerke-iot}
\todo[color=yellow]{IP-Adressen, Subnetze und Router vermitteln.}

\section{Modul 5: Smart-Home und Datenschutz}\label{sec:modul-datenschutz}
\todo[color=yellow]{Cloud, Lokalbetrieb und Datensouveränität vergleichen.}
